% SPDX-License-Identifier: CC-BY-4.0
% Copyright (c) 2026 Martin Schröder <info@swedishembedded.com>

\section{Results}\label{sec:results}

Every number in this section is produced from the raw records in \code{data/runs}
by \code{scripts/build_data.py}. The experiments are pilots. Of the four, E1 and E4 ran
as designed; E2 and E3 ran but did not produce a release, so what they measure is where
the loop stopped.

\subsection{E1: the funnel on a real corpus}\label{sec:e1}

The generator (Qwen3-4B) was asked for four kinds of task from the three pages of
Brain's documentation. The task report lists \mEOneTasks{} admitted tasks, covering
\mEOneParts{} of the \mEOnePartsTotal{} text parts of the corpus,
produced in \mEOneElapsedSecs{} seconds.\cl{M01} By kind,
recall admitted \mEOneAdmittedRecall, explain \mEOneAdmittedExplain, classify
\mEOneAdmittedClassify{} and the code-built denoise \mEOneAdmittedDenoise.\cl{M02}
The generator's proposals were rejected \mEOneRejectedRecall, \mEOneRejectedExplain,
\mEOneRejectedClassify{} and \mEOneRejectedDenoise{} times in the same order.\cl{M03}
\Cref{tab:e1} gives the rejections by reason. Of \mEOneRejectedTotal{} rejections,
\mEOneReasonNoSubject{} carry one reason, \code{no\_subject}: the generator did not
name what the task is about, which the admission rule of \cref{sec:admission}
requires. One was a duplicate and one was ungrounded.\cl{M04} We did not test why the generator omits the subject, nor whether a different brief
would admit more of them; this is a measurement of one generator on one corpus, not
a property of the admission rule.

\begin{table}[t]
\centering
\input{generated/e1_table}
\caption{E1: admissions and rejections of the task funnel (Qwen3-4B generator, three
documentation pages, one run).}\label{tab:e1}
\end{table}

\subsection{E2 and E3: a turn of the loop on invented facts}\label{sec:e2}

The invented facts (24 per generation, seed recorded in \code{data/runs/e2/build.json})
were not answered by the unadapted policy: \mETwoBaseGZero{} of 24 facts of generation~0,
\mETwoBaseGOne{} of generation~1 and \mETwoBaseGTwo{} of generation~2 had at least three of
four probes right, and \mETwoBaseAnchor{} of 40 anchor questions were right.\cl{M05}
The first \code{learn} reached the training stage with a small dataset
(\cref{tab:learn}): \mETwoGZeroTasks{} tasks were admitted from the 24-section
handbook, none of the \mETwoGZeroAttempts{} closed-book attempts of the policy was
verified correct, the teacher's answer was verified for \mETwoGZeroTeacherVerified{}
tasks, and the dataset held \mETwoGZeroDatasetRecords{} records.\cl{M06}
The second generation's \code{learn} ended with \mETwoGOneDatasetRecords{} records.\cl{M07}
Both commands exited with status \mETwoExitLearnGZero{} before a candidate existed,
with the message ``\mETwoGZeroRefusal''.\cl{M08} We have not diagnosed this refusal; it
occurred with eight and with four training records, the smallest datasets this
pipeline was given in any experiment here, and we do not know whether larger
datasets avoid it.

\begin{table}[t]
\centering
\input{generated/learn_table}
\caption{E2 and E3: stage counts of the two \code{learn} commands, read from the
command output kept in \code{data/runs/e2/raw/commands.log}.}\label{tab:learn}
\end{table}

Consequently the gate was never reached, no release was made and no alias moved. The
live policy answered the same probes as the base after each command
(\cref{tab:e2}): \mETwoAfterGZeroGZero{} of the first generation's facts after the first
command and \mETwoAfterGOneGOne{} of the second generation's after the second, and
\mETwoAfterGOneAnchor{} anchor questions, all as at the start; the release list was
empty (\mETwoReleasesFinal{} releases).\cl{M09} The rollback was refused with status
\mETwoExitRollback{}: ``\mETwoRollbackMessage''.\cl{M10} We therefore have \emph{no}
measurement of an improvement, of retention across generations, of a gate decision or of
a served release in a full turn: those parts of E2 and E3 are not measured. Two
qualified observations follow from what did run. The closed-book control worked as
designed: facts the policy cannot know were not answered by it (\mETwoBaseGZero{} of 24).\cl{M05}
The refusal to roll back an alias that points at no release is the specified behaviour
of that command.\cl{X06}

\begin{table}[t]
\centering
\input{generated/e2_table}
\caption{E2 and E3: facts answered (at least three of four probes) with Wilson 95\%
intervals, and anchor questions right, after each step. 24 facts per generation; the
policy did not change.}\label{tab:e2}
\end{table}

\subsection{E4: training the 27B model on one card}\label{sec:e4}

Brain's qualification at commit \mEFourCommit{} ran the 27B architecture on one 95~GiB
card; \cref{tab:e4} lists the runs. Probes use a randomly initialised model of the real
width and depth (64 layers), recomputation, rank-16 LoRA and one step, so their loss is not
a quality number. With an NF4 base, 8192 tokens fit in
\mEFourProbeNfFourEightOneNineTwoPeakMib{}~MiB at
\mEFourProbeNfFourEightOneNineTwoTokensPerS{} tokens/s;\cl{M11} with a bf16 base,
8192 tokens took \mEFourProbeBfOneSixEightOneNineTwoPeakMib{}~MiB and 16384 tokens
\mEFourProbeBfOneSixOneSixThreeEightFourPeakMib{}~MiB, both at
\mEFourProbeBfOneSixOneSixThreeEightFourTokensPerS{} tokens/s.\cl{M12,M13}
Through the command line the real FP8 checkpoint trained one step with an NF4 base and
with a bf16 base with a host peak of \mEFourCliNfFourOnestepHostRssMib{} and \mEFourCliBfOneSixOnestepHostRssMib{}~MiB respectively\cl{M14,M15} and, with an NF4 base,
40 steps drove the loss on eight training examples from
\mEFourCliNfFourFourZerostepLossFirst{} to \mEFourCliNfFourFourZerostepLossLast{}.\cl{M16}
That is a statement about fitting the training examples, not about generalisation: the
validation examples were two of the training examples. The command line's default
batch ran the device out of memory at 1024-token rows.\cl{M17}

\begin{table}[t]
\centering
\small
\input{generated/e4_table}
\caption{E4: qualification runs of the 27B architecture on one card (Brain's
qualification logs, copied unmodified apart from path names into \code{data/runs/qual}).}\label{tab:e4}
\end{table}

What E4 does not show: that a trained 27B adapter changes what the model generates (the
qualification found no generation path for an adapter on this model, and we do not claim
one), and any run beyond 16384 tokens or any NF4 run at 16384 on this commit.

\subsection{A documentation defect}\label{sec:defect}

Checking the release statements against the code (\cref{sec:gate}) found that the command-line
README placed a release under a directory that no code creates; the manifest is kept in the
experience database, the adapter in the content-addressed artifact store and the only
per-alias file is the serving manifest.\cl{X08} The README has been corrected in the same
change as this paper.
